\documentclass{article}

\usepackage{url}
\usepackage{amsmath,amssymb}
\usepackage{graphicx,svg}

\begin{document}
\title{Lecture 24\\State Space Techniques in CT}
\author{C.L. Wyatt}
\date{\today}
\maketitle

\section{Introduction}

Thus far we have only considered single-input, single-output (SISO) CT systems. These can be described by an N-th order LCCDE, inpulse response $h(t)$, or transfer function $H(s)$, and, if stable, via the frequency response $H(j\omega)$. In more complex systems we may be interested in multiple inputs, multiple outputs, and signals internal to the system.

Another term for signal is \textit{state} and we actually have had multiple states in our systems, however we have removed them for analysis by reducing the system description to a SISO inpulse response or transfer function. To see this, consider a second-order LCCDE

\[
a\, \frac{d^2y}{dt^2}(t) + b\,\frac{dy}{dt}(t) + c\, y(t) = x(t)
\]

Recall from your differential equations course that this can also be represented by two, first-order LCCDEs. Let $u_1(t) = y(t)$ and $u_2(t) = \frac{dy}{dt} = \frac{du_1(t)}{dt}$. Then $\frac{du_2}{dt} = \frac{d^2y}{dt^2}$ and the LCCDE is

\[
a\, \frac{du_2}{dt^2}(t) + b\,u_2(t) + c\, u_1(t) = x(t)
\]

Solving for $\frac{du_2}{dt}$

\[
\frac{du_2}{dt} = -\frac{b}{a}\, u_2(t) - \frac{c}{a}\, u_1(t) + \frac{1}{a}\, x(t)
\]

Noting $\frac{du_1}{dt} = u_2$ gives the system of ODEs. Switching to dot notation for derivatives as is tradition (since we will only be dealing with first derivatives), we obtain the system of equations:

\begin{align*}
  \dot{u}_1 (t) &= u_2(t) \\
  \dot{u}_2 (t) &= -\frac{b}{a}\, u_2(t) - \frac{c}{a}\, u_1(t) + \frac{1}{a}\, x(t)
\end{align*}

Collecting $u_1$ and $u_2$ into a vector

\[
\left[ \begin{array}{c}
  u_1 \\
  u_2
\end{array}\right]
\]

we can write the system as

\begin{align*}
\dot{u} &= \left[ \begin{array}{cc}
  0 & 1 \\
  -\frac{c}{a} & -\frac{b}{a}
\end{array}\right]\, u + \left[ \begin{array}{c}
  0 \\
  \frac{1}{a}
  \end{array}\right]\, x\\
&= A\, u + B\, x
\end{align*}
where $A$ is the 2x2 matrix and $B$ is a 2x1 matrix. This form, relating first derivatives of the state to a linear combinations of the states and inputs, is called a \textit{state equation}. Solving the state equation gives is the state/signal for both $u_1$ and $u_2$. To get the same result as solving the original LCCDE, we need the \textit{observation equation}, in this case a single output

\[
y = [1 \; 0]\, u = C\, u
\]

where $C$ is a 1x2 matrix.

To see where these internal states have been hiding consider the direct form I implementation of this system

\[
a\int\int y + b\int y + cy = x
\]

solving for $y$

\[
y = -\frac{a}{c}\int\int y - \frac{b}{c}\int y + \frac{1}{c}\, x
\]

Converting to the Laplace domain (assuming zero initial conditions

\[
Y(s) = -\frac{a}{c}\frac{1}{s^2} Y(s) - \frac{b}{c}\frac{1}{s} Y(s) + \frac{1}{c}\, X(s)
\]

As a block diagram this looks like

\begin{figure}[h!]
  \centering
  \includegraphics[scale=1, alt={block diagram for the example. see text.}]{figures/fig24_1}
  \caption{Block diagram for the example system.}
\end{figure}

The output of the integrator blocks correspond to the internal state variables.

TODO RESTART HERE


\end{document}
